% Options for packages loaded elsewhere
\PassOptionsToPackage{unicode}{hyperref}
\PassOptionsToPackage{hyphens}{url}
\PassOptionsToPackage{dvipsnames,svgnames,x11names}{xcolor}
\documentclass[
]{article}
\usepackage{xcolor}
\usepackage[a4paper, hmargin=2.2cm, vmargin=2.4cm, headheight=14.5pt, footskip=0.9cm, includehead, includefoot]{geometry}
\usepackage{amsmath,amssymb}
\setcounter{secnumdepth}{-\maxdimen} % remove section numbering
\usepackage{iftex}
\ifPDFTeX
  \usepackage[T1]{fontenc}
  \usepackage[utf8]{inputenc}
  \usepackage{textcomp} % provide euro and other symbols
\else % if luatex or xetex
  \usepackage{unicode-math} % this also loads fontspec
  \defaultfontfeatures{Scale=MatchLowercase}
  \defaultfontfeatures[\rmfamily]{Ligatures=TeX,Scale=1}
\fi
\usepackage{lmodern}
\ifPDFTeX\else
  % xetex/luatex font selection
\fi
% Use upquote if available, for straight quotes in verbatim environments
\IfFileExists{upquote.sty}{\usepackage{upquote}}{}
\IfFileExists{microtype.sty}{% use microtype if available
  \usepackage[]{microtype}
  \UseMicrotypeSet[protrusion]{basicmath} % disable protrusion for tt fonts
}{}
\makeatletter
\@ifundefined{KOMAClassName}{% if non-KOMA class
  \IfFileExists{parskip.sty}{%
    \usepackage{parskip}
  }{% else
    \setlength{\parindent}{0pt}
    \setlength{\parskip}{6pt plus 2pt minus 1pt}}
}{% if KOMA class
  \KOMAoptions{parskip=half}}
\makeatother
\usepackage{longtable,booktabs,array}
\newcounter{none} % for unnumbered tables
\usepackage{calc} % for calculating minipage widths
% Correct order of tables after \paragraph or \subparagraph
\usepackage{etoolbox}
\makeatletter
\patchcmd\longtable{\par}{\if@noskipsec\mbox{}\fi\par}{}{}
\makeatother
% Allow footnotes in longtable head/foot
\IfFileExists{footnotehyper.sty}{\usepackage{footnotehyper}}{\usepackage{footnote}}
\makesavenoteenv{longtable}
\setlength{\emergencystretch}{3em} % prevent overfull lines
\providecommand{\tightlist}{%
  \setlength{\itemsep}{0pt}\setlength{\parskip}{0pt}}
% ============================================================
% 土力学速查手册 —— 学术黑白 A4 版式（xelatex / pandoc）
% 主题：严谨、黑白、无彩色、无装饰
% ============================================================
\usepackage[fontset=windows]{ctex}            % 中文排版（宋体正文/黑体标题，自动伪粗体）
% 把常用符号划入 CJK 字体（SimSun），避免落到拉丁字体缺字：
% 带圈数字 ①-⑩、几何符号 ◆■□、箭头 →、希腊字母 αφ、数学符号 ≈
\xeCJKDeclareCharClass{CJK}{"2460->"24FF}
\xeCJKDeclareCharClass{CJK}{"25A0->"25FF}
\xeCJKDeclareCharClass{CJK}{"2190->"21FF}
\xeCJKDeclareCharClass{CJK}{"0370->"03FF}
\xeCJKDeclareCharClass{CJK}{"2248->"2248}% ≈（SimSun 有；− 等保留拉丁字体）
\usepackage{amsmath, amssymb}                 % 数学符号
\usepackage{booktabs}                         % 三线表
\usepackage{longtable}                        % 跨页表格
\usepackage{array}
\usepackage{fancyhdr}
\usepackage{enumitem}
\setlist{itemsep=2pt, topsep=4pt, leftmargin=1.6em}

% ---------- 页眉页脚：简洁学术 ----------
\pagestyle{fancy}
\fancyhf{}
\fancyhead[C]{\small\kaishu 土力学速查手册}
\fancyfoot[C]{\small\thepage}
\renewcommand{\headrulewidth}{0.4pt}
\renewcommand{\footrulewidth}{0pt}
\fancypagestyle{plain}{
  \fancyhf{}
  \fancyfoot[C]{\small\thepage}
  \renewcommand{\headrulewidth}{0pt}
}

% ---------- 表格：黑白三线风格 ----------
\renewcommand{\arraystretch}{1.35}
\setlength{\aboverulesep}{2pt}
\setlength{\belowrulesep}{2pt}

% ---------- 其他 ----------
\raggedbottom
\usepackage{bookmark}
\IfFileExists{xurl.sty}{\usepackage{xurl}}{} % add URL line breaks if available
\urlstyle{same}
% fallback for those not using the hyperref driver hyperxmp:
\makeatletter
\@ifundefined{xmpquote}{\newcommand{\xmpquote}[1]{#1}}{}
\makeatother
\hypersetup{
  colorlinks=true,
  linkcolor={black},
  filecolor={Maroon},
  citecolor={black},
  urlcolor={black},
  pdfcreator={LaTeX via pandoc}}

\author{}
\date{}

\begin{document}

\section{全书公式速查总表}\label{ux5168ux4e66ux516cux5f0fux901fux67e5ux603bux8868}

\begin{quote}
\textbf{用途}：考前最后冲刺 2
小时只翻这一份。按「主题」分栏，每个公式一行，附「适用/注意」与「出处（章节公式号）」便于回翻正文。详尽的参数、单位、推导与易错辨析见各章正文（目录见
README）。
\end{quote}

\begin{center}\rule{0.5\linewidth}{0.5pt}\end{center}

\subsection{〇、符号约定速查（全局统一）}\label{ux7b26ux53f7ux7ea6ux5b9aux901fux67e5ux5168ux5c40ux7edfux4e00}

重力加速度取 \(g=10\,\mathrm{m/s^2}\)，水重度
\(\gamma_w=10\,\mathrm{kN/m^3}\)，\(\rho_w=1\,\mathrm{g/cm^3}\)；含
\(\pi\) 的公式中三角函数一律用\textbf{弧度}。

{\def\LTcaptype{none} % do not increment counter
\begin{longtable}[]{@{}
  >{\raggedright\arraybackslash}p{(\linewidth - 4\tabcolsep) * \real{0.3333}}
  >{\raggedright\arraybackslash}p{(\linewidth - 4\tabcolsep) * \real{0.3333}}
  >{\raggedright\arraybackslash}p{(\linewidth - 4\tabcolsep) * \real{0.3333}}@{}}
\toprule\noalign{}
\begin{minipage}[b]{\linewidth}\raggedright
符号
\end{minipage} & \begin{minipage}[b]{\linewidth}\raggedright
含义
\end{minipage} & \begin{minipage}[b]{\linewidth}\raggedright
单位
\end{minipage} \\
\midrule\noalign{}
\endhead
\bottomrule\noalign{}
\endlastfoot
\(d_s\) & 土粒相对密度（比重） & --- \\
\(\omega,\ \omega_L,\ \omega_P\) & 含水率、液限、塑限 & \% \\
\(I_p,\ I_L\) & 塑性指数、液性指数 & ---（\(I_p\) 数值不带 \%） \\
\(e,\ n,\ S_r\) & 孔隙比、孔隙率、饱和度 & --- \\
\(\rho,\rho_d,\rho_{sat},\rho'\) & 密度、干密度、饱和密度、浮密度 &
\(\mathrm{g/cm^3}\) \\
\(\gamma,\gamma_d,\gamma_{sat},\gamma'\) &
重度、干重度、饱和重度、浮重度 & \(\mathrm{kN/m^3}\) \\
\(\sigma,\sigma',u\) & 总应力、有效应力、孔隙水压力 &
\(\mathrm{kPa}\) \\
\(\sigma_1,\sigma_3\) & 大、小主应力 & \(\mathrm{kPa}\) \\
\(c,\varphi\)；\(c',\varphi'\) & 总应力指标；有效应力指标 &
\(\mathrm{kPa}\)，° \\
\(\tau_f\) & 抗剪强度 & \(\mathrm{kPa}\) \\
\(c_u\) & 不排水抗剪强度（三轴 UU / 十字板） & \(\mathrm{kPa}\) \\
\(C_u\) & 不均匀系数（级配，第一章） & --- \\
\(C_c\) & ◆ 双义：第一章=曲率系数；第四章=压缩指数 & --- \\
\(C_s\) & 回弹指数（第四章） & --- \\
\(d_{10},d_{30},d_{60}\) & 级配曲线上累计含量对应粒径 & mm \\
\(k,\ i,\ j\) & 渗透系数、水力梯度、渗透力 &
\(\mathrm{cm/s}\)、---、\(\mathrm{kN/m^3}\) \\
\(E_s,\ a,\ m_v\) & 压缩模量、压缩系数、体积压缩系数 &
\(\mathrm{MPa}\)、\(\mathrm{MPa^{-1}}\) \\
\(p_c,\ p_0,\ OCR\) & 先期固结压力、现有自重应力、超固结比 &
\(\mathrm{kPa}\)、--- \\
\(c_v,\ T_v,\ U,\ H\) & 固结系数、时间因子、固结度、排水距离 &
\(\mathrm{m^2/a}\)、---、---、m \\
\(K_0,\ K_a,\ K_p\) & 静止、主动、被动土压力系数 & --- \\
\(E_a,E_0,E_p\) & 主动、静止、被动土压力合力（每米墙长） &
\(\mathrm{kN/m}\) \\
\(p,\ p_0,\ p_{cr},\ p_{1/4},\ p_u\) &
基底压力、附加压力、临塑荷载、临界荷载、极限承载力 & \(\mathrm{kPa}\) \\
\(f_{ak},\ f_a\) & 地基承载力特征值、深宽修正后特征值 &
\(\mathrm{kPa}\) \\
\(F_s\) & 安全系数 & --- \\
\(B,\ b,\ d,\ H\) & 基础宽度、条宽、埋深、挡墙/土层高 & m \\
\end{longtable}
}

\begin{center}\rule{0.5\linewidth}{0.5pt}\end{center}

\subsection{一、土的物理性质与三相换算（第一章）}\label{ux4e00ux571fux7684ux7269ux7406ux6027ux8d28ux4e0eux4e09ux76f8ux6362ux7b97ux7b2cux4e00ux7ae0}

{\def\LTcaptype{none} % do not increment counter
\begin{longtable}[]{@{}
  >{\raggedright\arraybackslash}p{(\linewidth - 6\tabcolsep) * \real{0.2424}}
  >{\raggedright\arraybackslash}p{(\linewidth - 6\tabcolsep) * \real{0.1818}}
  >{\raggedright\arraybackslash}p{(\linewidth - 6\tabcolsep) * \real{0.3939}}
  >{\raggedright\arraybackslash}p{(\linewidth - 6\tabcolsep) * \real{0.1818}}@{}}
\toprule\noalign{}
\begin{minipage}[b]{\linewidth}\raggedright
公式名
\end{minipage} & \begin{minipage}[b]{\linewidth}\raggedright
公式
\end{minipage} & \begin{minipage}[b]{\linewidth}\raggedright
适用 / 注意
\end{minipage} & \begin{minipage}[b]{\linewidth}\raggedright
出处
\end{minipage} \\
\midrule\noalign{}
\endhead
\bottomrule\noalign{}
\endlastfoot
含水率 & \(\omega=\dfrac{m_w}{m_s}\times100\%\) &
分母是\textbf{干土}质量 & 1-1 \\
密度/重度 & \(\rho=\dfrac{m}{V},\quad \gamma=\rho g\)（\(g=10\)） &
\(1\,\mathrm{g/cm^3}\leftrightarrow10\,\mathrm{kN/m^3}\) & 1-2 \\
土粒比重 & \(d_s=\dfrac{m_s}{V_s\rho_w}\) &
\(\rho_w=1\,\mathrm{g/cm^3}\) & 1-3 \\
孔隙比/孔隙率 & \(e=\dfrac{V_v}{V_s},\ n=\dfrac{V_v}{V}=\dfrac{e}{1+e}\)
& \(e=\dfrac{n}{1-n}\) & 1-4 \\
饱和度 & \(S_r=\dfrac{V_w}{V_v}\times100\%\) & 饱和土 \(S_r=1\) & 1-5 \\
干密度 & \(\rho_d=\dfrac{m_s}{V}\) & 压实控制指标 & 1-6 \\
饱和/浮密度 & \(\rho'=\rho_{sat}-\rho_w\) & 水下自重用
\(\gamma'=\gamma_{sat}-\gamma_w\) & 1-7 \\
\textbf{换算总式} &
\(e=\dfrac{d_s(1+\omega)}{\rho}-1,\ n=\dfrac{e}{1+e},\ S_r=\dfrac{\omega d_s}{e}\)
& 已知 \(\rho,\omega,d_s\) 全推；以 \(m_s=1\) 为基准可现场推导 &
1-8/1-9 \\
重度组 &
\(\gamma=\dfrac{d_s(1+\omega)}{1+e}\gamma_w,\ \gamma_d=\dfrac{d_s}{1+e}\gamma_w,\ \gamma_{sat}=\dfrac{d_s+e}{1+e}\gamma_w,\ \gamma'=\dfrac{d_s-1}{1+e}\gamma_w\)
& 给 \(d_s,e\) 直接出四个重度 & 1-10/1-11 \\
塑性指数 & \(I_p=\omega_L-\omega_P\) & 粉土 \(I_p\le10\) & 1-14 \\
液性指数 & \(I_L=\dfrac{\omega-\omega_P}{I_p}\) & \(I_L>1\) 流塑 &
1-15 \\
相对密实度 & \(D_r=\dfrac{e_{max}-e}{e_{max}-e_{min}}\) & 无黏性土 &
1-16 \\
级配评价 &
\(C_u=\dfrac{d_{60}}{d_{10}},\ C_c=\dfrac{d_{30}^2}{d_{60}d_{10}}\) &
良好：\(C_u\ge5\) 且 \(C_c=1\sim3\) & 1-17 \\
\end{longtable}
}

\begin{center}\rule{0.5\linewidth}{0.5pt}\end{center}

\subsection{二、渗透性与渗流（第二章）}\label{ux4e8cux6e17ux900fux6027ux4e0eux6e17ux6d41ux7b2cux4e8cux7ae0}

{\def\LTcaptype{none} % do not increment counter
\begin{longtable}[]{@{}
  >{\raggedright\arraybackslash}p{(\linewidth - 6\tabcolsep) * \real{0.2424}}
  >{\raggedright\arraybackslash}p{(\linewidth - 6\tabcolsep) * \real{0.1818}}
  >{\raggedright\arraybackslash}p{(\linewidth - 6\tabcolsep) * \real{0.3939}}
  >{\raggedright\arraybackslash}p{(\linewidth - 6\tabcolsep) * \real{0.1818}}@{}}
\toprule\noalign{}
\begin{minipage}[b]{\linewidth}\raggedright
公式名
\end{minipage} & \begin{minipage}[b]{\linewidth}\raggedright
公式
\end{minipage} & \begin{minipage}[b]{\linewidth}\raggedright
适用 / 注意
\end{minipage} & \begin{minipage}[b]{\linewidth}\raggedright
出处
\end{minipage} \\
\midrule\noalign{}
\endhead
\bottomrule\noalign{}
\endlastfoot
达西定律 & \(v=ki=k\dfrac{\Delta h}{L}\) & 层流；\(v\) 为虚拟流速 &
2-1 \\
实际流速 & \(v'=v/n\) & 孔隙中平均流速 & 2-2 \\
渗流量 & \(Q=kiA\) & --- & 2-3 \\
常水头试验 & \(k=\dfrac{VL}{Aht}\) & 砂土 & 2-4 \\
变水头试验 & \(k=2.3\dfrac{aL}{At}\lg\dfrac{h_1}{h_2}\) & 黏性土 &
2-5 \\
成层水平向 & \(k_x=\dfrac{\sum k_ih_i}{H}\) & 平行渗流 & 2-6 \\
成层竖直向 & \(k_z=\dfrac{H}{\sum h_i/k_i}\) & 垂直渗流，别倒置 & 2-6 \\
经验估算 & \(k\approx Cd_{10}^{\,2}\) & Hazen 式；无试验资料时估值 &
2-8 \\
渗透力 & \(j=i\gamma_w\) & 方向\textbf{顺水流}，体积力 & 2-9 \\
临界梯度 & \(i_{cr}=\dfrac{\gamma'}{\gamma_w}=\dfrac{d_s-1}{1+e}\) &
流土判据 \(i>i_{cr}\) & 2-11 \\
\end{longtable}
}

\begin{center}\rule{0.5\linewidth}{0.5pt}\end{center}

\subsection{三、土中应力（第三章）}\label{ux4e09ux571fux4e2dux5e94ux529bux7b2cux4e09ux7ae0}

{\def\LTcaptype{none} % do not increment counter
\begin{longtable}[]{@{}
  >{\raggedright\arraybackslash}p{(\linewidth - 6\tabcolsep) * \real{0.2424}}
  >{\raggedright\arraybackslash}p{(\linewidth - 6\tabcolsep) * \real{0.1818}}
  >{\raggedright\arraybackslash}p{(\linewidth - 6\tabcolsep) * \real{0.3939}}
  >{\raggedright\arraybackslash}p{(\linewidth - 6\tabcolsep) * \real{0.1818}}@{}}
\toprule\noalign{}
\begin{minipage}[b]{\linewidth}\raggedright
公式名
\end{minipage} & \begin{minipage}[b]{\linewidth}\raggedright
公式
\end{minipage} & \begin{minipage}[b]{\linewidth}\raggedright
适用 / 注意
\end{minipage} & \begin{minipage}[b]{\linewidth}\raggedright
出处
\end{minipage} \\
\midrule\noalign{}
\endhead
\bottomrule\noalign{}
\endlastfoot
自重应力 & \(\sigma_{cz}=\sum\gamma_ih_i\) & 水位以下用
\(\gamma'\)；隔水层顶面加 \(\gamma_w h_w\) & 3-1/3-2 \\
基底压力(中心) & \(p=\dfrac{N+G}{A}\) & \(G\) 含基础上的土重 & 3-5 \\
偏心（\(e\le B/6\)） &
\(p_{max,min}=\dfrac{N+G}{lB}\left(1\pm\dfrac{6e}{B}\right)\) &
核心区判别 & 3-6 \\
大偏心（\(e>B/6\)） & \(p_{max}=\dfrac{2(N+G)}{3al},\ a=\dfrac{B}{2}-e\)
& 三角形分布，\(p_{min}=0\) & 3-7 \\
基底附加压力 & \(p_0=p-\gamma_m d\) & 沉降计算\textbf{用它} & 3-8 \\
集中力解 & \(\sigma_z=\dfrac{3Pz^3}{2\pi R^5}\) & 布辛内斯克 & 3-9 \\
矩形角点法 & \(\sigma_z=\alpha_c p_0\) & 角点在 \(M\)
处，\(m=l/z,n=b/z\) & 3-10 \\
\textbf{平均附加应力系数} &
\(\bar\alpha_c=\dfrac{1}{z}\int_0^z\alpha_c\,dz\) &
应力面积法用（\textbf{平均}系数表，与角点 α 两套表） & 3-11 \\
条形荷载 &
\(\sigma_z=\dfrac{p_0}{\pi}\left[\beta+\sin\beta\cos(\alpha_1-\alpha_2)\right]\)
& 角度用弧度 & 3-13 \\
有效应力原理 & \(\sigma=\sigma'+u\) & 决定强度与变形的是 \(\sigma'\) &
3-14 \\
\end{longtable}
}

\begin{center}\rule{0.5\linewidth}{0.5pt}\end{center}

\subsection{四、压缩性、沉降与固结（第四章）}\label{ux56dbux538bux7f29ux6027ux6c89ux964dux4e0eux56faux7ed3ux7b2cux56dbux7ae0}

{\def\LTcaptype{none} % do not increment counter
\begin{longtable}[]{@{}
  >{\raggedright\arraybackslash}p{(\linewidth - 6\tabcolsep) * \real{0.2424}}
  >{\raggedright\arraybackslash}p{(\linewidth - 6\tabcolsep) * \real{0.1818}}
  >{\raggedright\arraybackslash}p{(\linewidth - 6\tabcolsep) * \real{0.3939}}
  >{\raggedright\arraybackslash}p{(\linewidth - 6\tabcolsep) * \real{0.1818}}@{}}
\toprule\noalign{}
\begin{minipage}[b]{\linewidth}\raggedright
公式名
\end{minipage} & \begin{minipage}[b]{\linewidth}\raggedright
公式
\end{minipage} & \begin{minipage}[b]{\linewidth}\raggedright
适用 / 注意
\end{minipage} & \begin{minipage}[b]{\linewidth}\raggedright
出处
\end{minipage} \\
\midrule\noalign{}
\endhead
\bottomrule\noalign{}
\endlastfoot
压缩系数 & \(a=\dfrac{e_1-e_2}{p_2-p_1}\) & 压缩性评价用 \(a_{1-2}\) &
4-1 \\
压缩模量 & \(E_s=\dfrac{1+e_1}{a}\) & 与 \(a\) 互推 & 4-2 \\
压缩指数 & \(C_c=\dfrac{\Delta e}{\lg(p_2/p_1)}\) & e-lgp 直线段 &
4-3 \\
超固结比 & \(OCR=\dfrac{p_c}{p_0}\) & \(OCR>1\) 超固结 & 4-4 \\
e-lgp 沉降 & 正常固结
\(s=\dfrac{C_c}{1+e_0}h\lg\dfrac{p_0+\Delta p}{p_0}\)；超固结分段（回弹用
\(C_s\)） & 按 OCR 与终点应力选式 & 4-8 \\
分层总和法 & \(s_i=\dfrac{e_1-e_2}{1+e_1}h_i\)（或
\(\dfrac{\Delta p_i}{E_s}h_i\)） & 压缩层底
\(\sigma_z\le0.2\sigma_{cz}\) & 4-5/4-6 \\
规范法 &
\(s=\psi_s\sum\dfrac{p_0}{E_s}\left(z_i\bar\alpha_i-z_{i-1}\bar\alpha_{i-1}\right)\)
& 平均附加应力系数 \(\bar\alpha\)（查表）+ \(\psi_s\) & 4-7 \\
固结系数 & \(c_v=\dfrac{kE_s}{\gamma_w}\) & 一维固结 & 4-9 \\
时间因子 & \(T_v=\dfrac{c_vt}{H^2}\) & 双面排水 \(H=H_0/2\) & 4-10 \\
固结度小 & \(U_t=2\sqrt{T_v/\pi}\) & \(0<U\le60\%\) & 4-12 \\
固结度大 & \(U_t=1-\dfrac{8}{\pi^2}e^{-\pi^2T_v/4}\) & \(U>60\%\) &
4-12 \\
\end{longtable}
}

\begin{quote}
注：4-12 的 \(U\)--\(T_v\)
关系对应\textbf{初始超孔压均匀分布}；三角形等非均匀分布不适用（见正文
4-12 旁说明）。
\end{quote}

\begin{center}\rule{0.5\linewidth}{0.5pt}\end{center}

\subsection{五、抗剪强度（第五章）}\label{ux4e94ux6297ux526aux5f3aux5ea6ux7b2cux4e94ux7ae0}

{\def\LTcaptype{none} % do not increment counter
\begin{longtable}[]{@{}
  >{\raggedright\arraybackslash}p{(\linewidth - 6\tabcolsep) * \real{0.2424}}
  >{\raggedright\arraybackslash}p{(\linewidth - 6\tabcolsep) * \real{0.1818}}
  >{\raggedright\arraybackslash}p{(\linewidth - 6\tabcolsep) * \real{0.3939}}
  >{\raggedright\arraybackslash}p{(\linewidth - 6\tabcolsep) * \real{0.1818}}@{}}
\toprule\noalign{}
\begin{minipage}[b]{\linewidth}\raggedright
公式名
\end{minipage} & \begin{minipage}[b]{\linewidth}\raggedright
公式
\end{minipage} & \begin{minipage}[b]{\linewidth}\raggedright
适用 / 注意
\end{minipage} & \begin{minipage}[b]{\linewidth}\raggedright
出处
\end{minipage} \\
\midrule\noalign{}
\endhead
\bottomrule\noalign{}
\endlastfoot
库仑公式(总应力) & \(\tau_f=c+\sigma\tan\varphi\) & 砂土 \(c=0\) &
5-1 \\
库仑公式(有效) & \(\tau_f=c'+(\sigma-u)\tan\varphi'\) & 真正决定强度的是
\(\sigma'\) & 5-2 \\
\textbf{任意平面应力} &
\(\sigma=\dfrac{\sigma_1+\sigma_3}{2}+\dfrac{\sigma_1-\sigma_3}{2}\cos2\theta,\ \tau=\dfrac{\sigma_1-\sigma_3}{2}\sin2\theta\)
& 莫尔圆几何基础 & 5-3 \\
破坏判定 & 莫尔圆与包线 \(\tau=c+\sigma\tan\varphi\) 相切 &
相切=极限平衡，越过=破坏 & 5-4 \\
极限平衡条件 &
\(\sigma_1=\sigma_3\tan^2\left(45^{\circ}+\dfrac{\varphi}{2}\right)+2c\tan\left(45^{\circ}+\dfrac{\varphi}{2}\right)\)
& 与 \(\sigma_3=f(\sigma_1)\) 两式等价 & 5-5 \\
破裂面方位 &
\(\alpha_f=45^{\circ}+\dfrac{\varphi}{2}\)（与大主应力作用面夹角） & ---
& 5-6 \\
\textbf{无侧限强度} & \(q_u=2c_u,\ c_u=\dfrac{q_u}{2}\) & \(\varphi=0\)
特例，饱和黏土 & 5-7 \\
孔压系数 &
\(\Delta u=B\left[\Delta\sigma_3+A(\Delta\sigma_1-\Delta\sigma_3)\right]\)
& 饱和土 \(B=1\) & 5-11 \\
\end{longtable}
}

\textbf{试验与指标对照（简答高频）}：直剪 快剪/固结快剪/慢剪 ≈ 三轴
UU/CU/CD；饱和黏土 UU 得 \(c_u,\varphi_u=0\)；CU 同时得总应力
\(c_{cu},\varphi_{cu}\) 与有效应力 \(c',\varphi'\)；慢剪 ≈
有效应力指标。

\begin{center}\rule{0.5\linewidth}{0.5pt}\end{center}

\subsection{六、土压力与挡土墙（第六章）}\label{ux516dux571fux538bux529bux4e0eux6321ux571fux5899ux7b2cux516dux7ae0}

{\def\LTcaptype{none} % do not increment counter
\begin{longtable}[]{@{}
  >{\raggedright\arraybackslash}p{(\linewidth - 6\tabcolsep) * \real{0.2424}}
  >{\raggedright\arraybackslash}p{(\linewidth - 6\tabcolsep) * \real{0.1818}}
  >{\raggedright\arraybackslash}p{(\linewidth - 6\tabcolsep) * \real{0.3939}}
  >{\raggedright\arraybackslash}p{(\linewidth - 6\tabcolsep) * \real{0.1818}}@{}}
\toprule\noalign{}
\begin{minipage}[b]{\linewidth}\raggedright
公式名
\end{minipage} & \begin{minipage}[b]{\linewidth}\raggedright
公式
\end{minipage} & \begin{minipage}[b]{\linewidth}\raggedright
适用 / 注意
\end{minipage} & \begin{minipage}[b]{\linewidth}\raggedright
出处
\end{minipage} \\
\midrule\noalign{}
\endhead
\bottomrule\noalign{}
\endlastfoot
三者关系 & \(E_a<E_0<E_p\) & 墙远离/静止/挤压 & 6-1 \\
静止系数 & \(K_0=1-\sin\varphi'\) & Jaky 正常固结土 & 6-2 \\
朗肯系数 &
\(K_a=\tan^2\left(45^{\circ}-\dfrac{\varphi}{2}\right),\ K_p=\tan^2\left(45^{\circ}+\dfrac{\varphi}{2}\right)\)
& \(K_aK_p=1\) & 6-3 \\
主动(无黏) & \(p_a=\gamma zK_a,\ E_a=\dfrac{1}{2}\gamma H^2K_a\) &
三角形，作用点距底 \(H/3\) & 6-4 \\
主动(黏性) &
\(p_a=\gamma zK_a-2c\sqrt{K_a},\ z_0=\dfrac{2c}{\gamma\sqrt{K_a}},\ E_a=\dfrac{1}{2}\gamma(H-z_0)^2K_a\)
& 忽略墙顶拉应力区 & 6-5 \\
被动(无黏) & \(E_p=\dfrac{1}{2}\gamma H^2K_p\) & --- & 6-6 \\
被动(黏性) & \(E_p=\dfrac{1}{2}\gamma H^2K_p+2cH\sqrt{K_p}\) & 梯形分布
& 6-7 \\
抗倾覆 & \(K_t=\dfrac{Wx_0}{E_ay_0}\ge1.5\) & 对墙趾取矩 & 6-11 \\
抗滑 & \(K_s=\dfrac{(W+E_{ay})\mu}{E_{ax}}\ge1.3\) & 用水平分量 &
6-12 \\
\end{longtable}
}

\textbf{库仑适用}：墙背倾斜、粗糙（\(\delta\ne0\)）、填土面倾斜
\(\beta\)；\(K_a\) 一般式见第六章公式 6-9；当
\(\alpha=90^{\circ},\delta=0,\beta=0\) 退化为朗肯。

\begin{center}\rule{0.5\linewidth}{0.5pt}\end{center}

\subsection{七、土坡稳定（第七章）}\label{ux4e03ux571fux5761ux7a33ux5b9aux7b2cux4e03ux7ae0}

{\def\LTcaptype{none} % do not increment counter
\begin{longtable}[]{@{}
  >{\raggedright\arraybackslash}p{(\linewidth - 6\tabcolsep) * \real{0.2424}}
  >{\raggedright\arraybackslash}p{(\linewidth - 6\tabcolsep) * \real{0.1818}}
  >{\raggedright\arraybackslash}p{(\linewidth - 6\tabcolsep) * \real{0.3939}}
  >{\raggedright\arraybackslash}p{(\linewidth - 6\tabcolsep) * \real{0.1818}}@{}}
\toprule\noalign{}
\begin{minipage}[b]{\linewidth}\raggedright
公式名
\end{minipage} & \begin{minipage}[b]{\linewidth}\raggedright
公式
\end{minipage} & \begin{minipage}[b]{\linewidth}\raggedright
适用 / 注意
\end{minipage} & \begin{minipage}[b]{\linewidth}\raggedright
出处
\end{minipage} \\
\midrule\noalign{}
\endhead
\bottomrule\noalign{}
\endlastfoot
安全系数定义 & \(F_s=\dfrac{\text{抗滑}}{\text{滑动力(矩)}}\) &
\(F_s\ge1.0\) & 7-1 \\
无黏性土坡 & \(F_s=\dfrac{\tan\varphi}{\tan\beta}\) & 与坡高无关 &
7-2 \\
\textbf{顺坡渗流} &
\(F_s=\dfrac{\gamma'}{\gamma_{sat}}\cdot\dfrac{\tan\varphi}{\tan\beta}\)
& 无黏性土坡渗流平行坡面，比 7-2 多乘 \(\gamma'/\gamma_{sat}\) & 7-7 \\
整体圆弧(φ=0) & \(F_s=\dfrac{c_uLR}{Wd}\) & 均质饱和软黏土 & 7-3 \\
稳定数 & \(N_s=\dfrac{c}{\gamma H_{cr}}\) & 泰勒图解法 & 7-4 \\
瑞典条分法 &
\(F_s=\dfrac{\sum\left[c_il_i+(W_i\cos\alpha_i-u_il_i)\tan\varphi_i\right]}{\sum W_i\sin\alpha_i}\)
& 忽略条间力，偏保守 & 7-5 \\
简化毕肖普 &
\(F_s=\dfrac{\sum\dfrac{c'b+(W-ub)\tan\varphi'}{\cos\alpha(1+\tan\alpha\tan\varphi'/F_s)}}{\sum W\sin\alpha}\)
& 有效应力、隐式迭代 & 7-6 \\
\end{longtable}
}

\begin{center}\rule{0.5\linewidth}{0.5pt}\end{center}

\subsection{八、地基承载力（第八章）}\label{ux516bux5730ux57faux627fux8f7dux529bux7b2cux516bux7ae0}

{\def\LTcaptype{none} % do not increment counter
\begin{longtable}[]{@{}
  >{\raggedright\arraybackslash}p{(\linewidth - 6\tabcolsep) * \real{0.2424}}
  >{\raggedright\arraybackslash}p{(\linewidth - 6\tabcolsep) * \real{0.1818}}
  >{\raggedright\arraybackslash}p{(\linewidth - 6\tabcolsep) * \real{0.3939}}
  >{\raggedright\arraybackslash}p{(\linewidth - 6\tabcolsep) * \real{0.1818}}@{}}
\toprule\noalign{}
\begin{minipage}[b]{\linewidth}\raggedright
公式名
\end{minipage} & \begin{minipage}[b]{\linewidth}\raggedright
公式
\end{minipage} & \begin{minipage}[b]{\linewidth}\raggedright
适用 / 注意
\end{minipage} & \begin{minipage}[b]{\linewidth}\raggedright
出处
\end{minipage} \\
\midrule\noalign{}
\endhead
\bottomrule\noalign{}
\endlastfoot
临塑荷载 &
\(p_{cr}=\dfrac{\pi\left(\gamma_1d+c\cot\varphi\right)}{\cot\varphi-\pi/2+\varphi}+\gamma_1d\)
& 塑性区深度=0；角度用弧度 & 8-1 \\
临界荷载 &
\(p_{1/4}=p_{cr}+\dfrac{\pi\gamma B}{4\left(\cot\varphi-\pi/2+\varphi\right)}\)
& 塑性区深 \(B/4\)；φ=0 时 B 项消失 & 8-2 \\
太沙基(条形) & \(p_u=cN_c+\gamma_1dN_q+\dfrac{1}{2}\gamma BN_\gamma\) &
「c+埋深+宽度」三项 & 8-3 \\
方形 & \(p_u=1.3cN_c+\gamma_1dN_q+0.4\gamma BN_\gamma\) & 形状系数别漏 &
8-4 \\
圆形 & \(p_u=1.3cN_c+\gamma_1dN_q+0.6\gamma RN_\gamma\) & --- & 8-4 \\
承载力系数 &
\(N_q=\dfrac{e^{(3\pi/2-\varphi)\tan\varphi}}{2\cos^2(45^{\circ}+\varphi/2)},\ N_c=(N_q-1)\cot\varphi\)
& 查表为主 & 8-5 \\
\textbf{深宽修正} &
\(f_a=f_{ak}+\eta_b\gamma(b-3)+\eta_d\gamma_m(d-0.5)\) &
\(b\le6,\ d\ge0.5\)；\(\gamma\) 在下、\(\gamma_m\) 在上 & 8-6 \\
\end{longtable}
}

\begin{center}\rule{0.5\linewidth}{0.5pt}\end{center}

\subsection{九、高频易混对（考前 30
秒自查）}\label{ux4e5dux9ad8ux9891ux6613ux6df7ux5bf9ux8003ux524d-30-ux79d2ux81eaux67e5}

{\def\LTcaptype{none} % do not increment counter
\begin{longtable}[]{@{}
  >{\raggedright\arraybackslash}p{(\linewidth - 2\tabcolsep) * \real{0.4706}}
  >{\raggedright\arraybackslash}p{(\linewidth - 2\tabcolsep) * \real{0.5294}}@{}}
\toprule\noalign{}
\begin{minipage}[b]{\linewidth}\raggedright
易混对
\end{minipage} & \begin{minipage}[b]{\linewidth}\raggedright
判别记法
\end{minipage} \\
\midrule\noalign{}
\endhead
\bottomrule\noalign{}
\endlastfoot
\(\rho_d\) vs \(d_s\) & 干密度=整体去除水；比重=纯土粒相对水 \\
\(n\) vs \(e\) & \(n\) 分母 ×总 V；\(e\) 分母 ×\(V_s\) \\
\(C_c\)（曲率系数 vs 压缩指数） & 第一章级配评价 vs 第四章
e-lgp------相同符号、含义不同 \\
\(C_u\) vs \(c_u\) & 不均匀系数（第一章） vs
不排水强度（第五章），看清大小写 \\
\(I_p\) vs \(I_L\) & 一个是范围（材料属性），一个是状态 \\
\(a\) vs \(E_s\) & \(E_s=\dfrac{1+e_1}{a}\)，反向 \\
\(\gamma\) vs \(\gamma'\) & 水位以下一律 \(\gamma'\) \\
\(p\) vs \(p_0\) & 沉降计算用 \(p_0=p-\gamma_md\) \\
总应力 vs 有效应力 \(c,\varphi\) & 用 \(\sigma'\) 得
\(c',\varphi'\)，取决于排水 \\
朗肯 vs 库仑 & 竖直/光滑/水平 → 朗肯；倾斜/粗糙/倾斜 → 库仑 \\
\(K_a\) vs \(K_p\) & \(K_p=1/K_a\) \\
\(E_s\) 沉降 vs \(U\) & 先 \(T_v=\dfrac{c_vt}{H^2}\)，再查 \(U-T_v\) \\
均匀 vs 三角形初始孔压 & 固结度 \(U\)--\(T_v\)
曲线不同，均匀分布式不可直接套非均匀分布 \\
临塑 vs 临界荷载 & \(p_{1/4}=p_{cr}+\)「宽度项」，φ=0 时二者相等 \\
\end{longtable}
}

\begin{center}\rule{0.5\linewidth}{0.5pt}\end{center}

\emph{各公式的参数说明、单位与推导细节请回翻对应章节。祝考试顺利。}

\end{document}
